\documentclass{ximera}

\input{../../../preamble.tex}


\definecolor{fvdnavy}{HTML}{071D43}
\definecolor{fvdcyan}{HTML}{20BFC6}
\definecolor{fvdcyanlight}{HTML}{EFFCFB}
\definecolor{fvdmagenta}{HTML}{F10D69}
\definecolor{fvdmagentalight}{HTML}{FFF1F7}
\definecolor{fvdtext}{HTML}{163044}





%=================================================
% Pedagogische omgevingen
% PDF: tcolorbox
% HTML: learning-box
%=================================================

\newenvironment{voorkennis}
{%
    \ifdefined\HCode
        \HCode{%
<section class="learning-box learning-box--prior">
<div class="learning-box__header">
<span class="learning-box__icon" aria-hidden="true"></span>
<span class="learning-box__title">Voorkennis</span>
</div>
<div class="learning-box__content">
        }%
        \begin{itemize}
    \else
       \begin{tcolorbox}[
    enhanced,
    breakable,
    colback=fvdcyanlight,
    colframe=fvdcyan,
    coltext=fvdtext,
    boxrule=0.5pt,
    leftrule=4pt,
    arc=4mm,
    outer arc=4mm,
    left=7mm,
    right=6mm,
    top=5mm,
    bottom=5mm,
    before skip=6mm,
    after skip=6mm,
    title=Voorkennis,
    coltitle=fvdcyan,
    fonttitle=\bfseries\Large,
    attach title to upper={\par\smallskip},
    boxed title style={
        empty,
        left=0mm,
        right=0mm,
        top=0mm,
        bottom=0mm
    },
    overlay={
        \draw[
            fvdcyan,
            line width=0.5pt
        ]
        ([xshift=42mm,yshift=-13mm]frame.north west)
        --
        ([xshift=-7mm,yshift=-13mm]frame.north east);
    }
]
        \begin{itemize}
    \fi
}
{%
    \end{itemize}
    \ifdefined\HCode
        \HCode{%
</div>
</section>
        }%
    \else
        \end{tcolorbox}
    \fi
}


\newenvironment{leerdoelen}
{%
    \ifdefined\HCode
        \HCode{%
<section class="learning-box learning-box--goals">
<div class="learning-box__header">
<span class="learning-box__icon" aria-hidden="true"></span>
<span class="learning-box__title">Leerdoelen</span>
</div>
<div class="learning-box__content">
        }%
        \begin{itemize}
    \else
\begin{tcolorbox}[
    enhanced,
    breakable,
    colback=fvdmagentalight,
    colframe=fvdmagenta,
    coltext=fvdtext,
    boxrule=0.5pt,
    leftrule=4pt,
    arc=4mm,
    outer arc=4mm,
    left=7mm,
    right=6mm,
    top=5mm,
    bottom=5mm,
    before skip=6mm,
    after skip=6mm,
    title=Leerdoelen,
    coltitle=fvdmagenta,
    fonttitle=\bfseries\Large,
    attach title to upper={\par\smallskip},
    boxed title style={
        empty,
        left=0mm,
        right=0mm,
        top=0mm,
        bottom=0mm
    },
    overlay={
        \draw[
            fvdmagenta,
            line width=0.5pt
        ]
        ([xshift=42mm,yshift=-13mm]frame.north west)
        --
        ([xshift=-7mm,yshift=-13mm]frame.north east);
    }
]
        \begin{itemize}
    \fi
}
{%
    \end{itemize}
    \ifdefined\HCode
        \HCode{%
</div>
</section>
        }%
    \else
        \end{tcolorbox}
    \fi
}


\addPrintStyle{../../..}

\title{In cirkels bewegen}
\author{Ingmar Herreman}

\begin{abstract}
Een voorwerp kan versnellen zonder sneller of trager te gaan. Dat klinkt
vreemd, maar precies dat gebeurt bij een eenparige cirkelbeweging:
de grootte van de snelheid blijft constant, terwijl de richting van de
snelheidsvector voortdurend verandert.

In dit hoofdstuk onderzoeken we periodieke bewegingen en bouwen we de
eenparige cirkelbeweging systematisch op. We verbinden periode en
frequentie met hoekpositie en hoeksnelheid en onderzoeken hoe de
baansnelheid afhangt van de straal en de hoeksnelheid.

Daarna staat de centrale vraag van het hoofdstuk centraal: waarom is er
bij een eenparige cirkelbeweging toch een versnelling? Via de verandering
van de snelheidsvector komen we tot de centripetale versnelling. We
onderzoeken de verbanden tussen straal, periode, hoeksnelheid,
baansnelheid en centripetale versnelling zowel theoretisch als
experimenteel.

De cirkelbeweging vormt zo een belangrijke schakel tussen kinematica en
dynamica. In het volgende hoofdstuk onderzoeken we welke resulterende
kracht nodig is om zo'n versnelling te veroorzaken.
\end{abstract}

\begin{document}

% FVD_CHAPTER_DATA_START
% Automatisch gegenereerd uit index.tex.
% Niet handmatig aanpassen.
\fvdchapterdata{1}{Beweging en kracht — Van beweging beschrijven tot energie begrijpen}{5}
% FVD_CHAPTER_DATA_END



\maketitle

% ============================================================
% 1 VOORKENNIS EN LEERDOELEN
% ============================================================

\begin{voorkennis}
  \item Je kan een beweging beschrijven met positie, snelheid en versnelling.
  \item Je weet dat snelheid en versnelling vectoriële grootheden zijn.
  \item Je kan vectoren tekenen, vergelijken en ontbinden.
  \item Je kent de omtrek van een cirkel: $2\pi r$.
  \item Je kan werken met graden en radialen.
  \item Je kan recht evenredige, omgekeerd evenredige en kwadratische
        verbanden herkennen en interpreteren.
\end{voorkennis}

\begin{leerdoelen}
  \item Je kan periode en frequentie van een periodieke beweging bepalen
        en interpreteren.
  \item Je kan hoekverplaatsing en hoeksnelheid gebruiken om een rotatie
        te beschrijven.
  \item Je kan het verband tussen $T$, $f$ en $\omega$ analyseren en
        kwantificeren.
  \item Je kan het onderscheid tussen hoeksnelheid en baansnelheid uitleggen.
  \item Je kan het verband $v=r\omega$ analyseren en toepassen.
  \item Je kan positie- en snelheidsvectoren bij een eenparige
        cirkelbeweging tekenen en interpreteren.
  \item Je kan verklaren waarom een eenparige cirkelbeweging een versnelde
        beweging is.
  \item Je kan de richting en grootte van de centripetale versnelling
        bepalen.
  \item Je kan de verbanden tussen $r$, $T$, $f$, $\omega$, $v$ en $a_c$
        combineren om nieuwe situaties te analyseren.
  \item Je kan experimentele gegevens over een cirkelbeweging verwerken,
        verbanden onderzoeken en conclusies formuleren.
\end{leerdoelen}


% ============================================================
% 2 OPENING
% ============================================================

\FVDsection{Versnellen zonder sneller te gaan?}

Een auto rijdt met constante snelheid door een bocht. Een satelliet
beweegt bijna met constante snelheid rond de aarde. Een punt op de rand
van een draaiende schijf legt elke seconde evenveel afstand af.

Zijn deze bewegingen onversneld?

Bij een rechtlijnige beweging zou een constante snelheid betekenen dat
de snelheidsvector constant is. In een cirkel verandert echter
voortdurend de \emph{richting} van de snelheid.

\begin{oefening}[Voorspel vóór je rekent]{\oefB\ESS}

Een fietser rijdt met een constante snelheidsgrootte van
$6,0\,\mathrm{m/s}$ door een cirkelvormige bocht.

\begin{question}
Is de grootte van zijn snelheid constant?
\end{question}

\begin{question}
Is zijn snelheidsvector constant? Motiveer.
\end{question}

\begin{question}
Verwacht je dat zijn versnelling nul is? Noteer eerst je intuïtieve
antwoord.
\end{question}

\begin{question}
In welke richting zou een eventuele versnelling volgens jou wijzen?
\end{question}

Bewaar je antwoorden. We keren er later naar terug.

\end{oefening}


% ============================================================
% 3 PERIODIEKE BEWEGING
% ============================================================

\FVDsection{Periodieke bewegingen}

Veel bewegingen herhalen zich: een wiel draait rond, een ventilator
roteert, de aarde draait om haar as en een punt op een draaimolen keert
steeds opnieuw naar dezelfde positie terug.

\begin{definitie}
Een \textbf{periodieke beweging} is een beweging die zich na gelijke
tijdsintervallen herhaalt.
\end{definitie}

Eén volledige herhaling noemen we een \textbf{cyclus}.


\FVDsubsection{Periode}

De \textbf{periode} $T$ is de tijdsduur van één volledige cyclus.

\[
\boxed{T=\frac{\Delta t}{N}}
\]

waarbij $N$ het aantal volledige cycli is.

De SI-eenheid is de seconde:

\[
[T]=\mathrm{s}.
\]


\FVDsubsection{Frequentie}

De \textbf{frequentie} $f$ is het aantal volledige cycli per seconde:

\[
\boxed{f=\frac{N}{\Delta t}}.
\]

De SI-eenheid is hertz:

\[
[f]=\mathrm{Hz}=\mathrm{s^{-1}}.
\]

Omdat periode en frequentie elkaars inverse zijn:

\[
\boxed{f=\frac1T}
\qquad\text{en}\qquad
\boxed{T=\frac1f}.
\]


\begin{voorbeeld}[Een draaiend wiel]

Een wiel maakt $45$ omwentelingen in $30,0\,\mathrm{s}$.

\[
f=\frac{45}{30,0}=1,50\,\mathrm{Hz}.
\]

De periode is:

\[
T=\frac1{1,50}=0,667\,\mathrm{s}.
\]

Elke omwenteling duurt dus ongeveer $0,667\,\mathrm{s}$.

\end{voorbeeld}


\begin{oefening}[Periode en frequentie]{\oefB\ESS}

Een laboratoriumrotor maakt $1800$ omwentelingen per minuut.

\begin{question}
Bepaal de frequentie in hertz.
\end{question}

\begin{question}
Bepaal de periode.
\end{question}

\begin{question}
Hoeveel omwentelingen maakt de rotor in $2,50\,\mathrm{s}$?
\end{question}

\end{oefening}


% ============================================================
% 4 EENPARIGE CIRKELBEWEGING
% ============================================================

\FVDsection{De eenparige cirkelbeweging}

\begin{definitie}
Bij een \textbf{eenparige cirkelbeweging} (ECB) beweegt een punt langs
een cirkelbaan met een constante \emph{grootte} van de baansnelheid.
\end{definitie}

Het woord \emph{eenparig} betekent hier dus niet dat de
snelheidsvector constant is.

\[
|\vec v|=\text{constant},
\]

maar de richting van $\vec v$ verandert voortdurend.

Dat onderscheid wordt de kern van dit hoofdstuk.


\FVDsubsection{Hoekpositie}

De positie op de cirkel kunnen we beschrijven met een hoek $\theta$.

Voor een booglengte $s$ op een cirkel met straal $r$ geldt, als
$\theta$ in radialen wordt uitgedrukt:

\[
\boxed{s=r\theta}.
\]

Eén volledige omwenteling komt overeen met:

\[
\boxed{2\pi\ \mathrm{rad}=360^\circ}.
\]


\FVDsubsection{Hoeksnelheid}

De gemiddelde hoeksnelheid is:

\[
\omega_{\mathrm{gem}}
=
\frac{\Delta\theta}{\Delta t}.
\]

Bij een ECB is de hoeksnelheid constant:

\[
\boxed{\omega=\frac{\Delta\theta}{\Delta t}}.
\]

De eenheid is:

\[
[\omega]=\mathrm{rad/s}.
\]

Tijdens één periode verandert de hoek met $2\pi$ rad. Daarom:

\[
\omega
=
\frac{2\pi}{T}.
\]

Met $f=1/T$ volgt ook:

\[
\boxed{
\omega=\frac{2\pi}{T}=2\pi f.
}
\]


\begin{oefening}[Van periode naar hoeksnelheid]{\oefB\ESS}

Een draaimolen maakt één omwenteling in $8,0\,\mathrm{s}$.

\begin{question}
Bepaal $T$.
\end{question}

\begin{question}
Bepaal $f$.
\end{question}

\begin{question}
Bepaal $\omega$.
\end{question}

\begin{question}
Welke hoek, in radialen en in graden, doorloopt de draaimolen in
$3,0\,\mathrm{s}$?
\end{question}

\end{oefening}


% ============================================================
% 5 BAANSNELHEID
% ============================================================

\FVDsection{Van hoeksnelheid naar baansnelheid}

Tijdens één volledige omwenteling legt een punt op afstand $r$ van het
middelpunt een afstand

\[
2\pi r
\]

af in een tijd $T$.

Daarom is de grootte van de baansnelheid:

\[
v=\frac{2\pi r}{T}.
\]

Omdat

\[
\omega=\frac{2\pi}{T},
\]

vinden we:

\[
\boxed{v=r\omega}.
\]

Dit is een fundamenteel verband.


\FVDsubsection{Zelfde hoeksnelheid, andere baansnelheid}

Denk aan twee punten op dezelfde stijve draaiende schijf.

Beide punten maken in dezelfde tijd dezelfde hoekverplaatsing. Ze
hebben dus dezelfde $\omega$, $T$ en $f$.

Het punt verder van het middelpunt legt in diezelfde tijd echter een
grotere booglengte af. Daarom is zijn baansnelheid groter.

Bij constante $\omega$:

\[
\boxed{v\sim r}.
\]


\begin{oefening}[Twee punten op één schijf]{\oefU\ESS}

Een schijf draait met constante hoeksnelheid. Punt A ligt op
$r_A=0,10\,\mathrm{m}$ van het middelpunt en punt B op
$r_B=0,30\,\mathrm{m}$.

\begin{question}
Vergelijk $T_A$ en $T_B$.
\end{question}

\begin{question}
Vergelijk $f_A$ en $f_B$.
\end{question}

\begin{question}
Vergelijk $\omega_A$ en $\omega_B$.
\end{question}

\begin{question}
Bepaal de verhouding $\dfrac{v_B}{v_A}$ zonder een concrete
hoeksnelheid te kiezen.
\end{question}

\begin{question}
Leg het resultaat uit vanuit de afgelegde booglengte.
\end{question}

\end{oefening}


\begin{voorbeeld}[Een punt op een draaiende schijf]

Een schijf draait met $f=2,50\,\mathrm{Hz}$. Een punt bevindt zich
$0,20\,\mathrm{m}$ van het middelpunt.

Eerst:

\[
\omega=2\pi f
=2\pi\cdot2,50
\approx15,7\,\mathrm{rad/s}.
\]

Daarna:

\[
v=r\omega
=0,20\cdot15,7
\approx3,14\,\mathrm{m/s}.
\]

\end{voorbeeld}


\begin{oefening}[CD-speler]{\oefB}

Een punt op een draaiende schijf bevindt zich op
$5,0\,\mathrm{cm}$ van het middelpunt. De schijf draait met
$300\,\mathrm{omw/min}$.

\begin{question}
Zet de rotatiesnelheid om naar hertz.
\end{question}

\begin{question}
Bepaal de hoeksnelheid.
\end{question}

\begin{question}
Bepaal de baansnelheid van het punt.
\end{question}

\end{oefening}


% ============================================================
% 6 POSITIEVECTOR EN SNELHEIDSVECTOR
% ============================================================

\FVDsection{De vectoriële beschrijving}

We kiezen het middelpunt van de cirkel als oorsprong.

De positievector $\vec r$ wijst van het middelpunt naar het bewegende
punt. Zijn grootte is constant:

\[
|\vec r|=r.
\]

Maar zijn richting verandert voortdurend.


\FVDsubsection{De snelheidsvector}

De snelheidsvector is op elk punt \textbf{raaklijnig aan de cirkel}.

Daarom geldt:

\[
\boxed{\vec v\perp\vec r}.
\]

Bij een ECB blijft:

\[
|\vec v|=v
\]

constant, maar $\vec v$ zelf niet.

\begin{opmerking}
Een constante \emph{snelheidsgrootte} is niet hetzelfde als een
constante \emph{snelheidsvector}. Een vector is pas constant wanneer
zowel zijn grootte als zijn richting constant zijn.
\end{opmerking}


\begin{oefening}[Vectoren tekenen]{\oefU\ESS}

Teken een cirkel en kies vier punten die telkens een kwart omwenteling
van elkaar verwijderd zijn.

\begin{question}
Teken in elk punt de positievector $\vec r$.
\end{question}

\begin{question}
Teken in elk punt de snelheidsvector $\vec v$ voor een beweging tegen
de wijzers van de klok in.
\end{question}

\begin{question}
Wat kan je zeggen over de grootte van de vier snelheidsvectoren?
\end{question}

\begin{question}
Wat kan je zeggen over hun richting?
\end{question}

\begin{question}
Waarom kan de versnelling dus niet nul zijn?
\end{question}

\end{oefening}


% ============================================================
% 7 WAAROM VERSNELLING?
% ============================================================

\FVDsection{Waarom is een ECB een versnelde beweging?}

Versnelling beschrijft de verandering van de snelheidsvector:

\[
\vec a
=
\frac{\Delta\vec v}{\Delta t}
\]

voor een gemiddeld tijdsinterval, en ogenblikkelijk als limiet voor
$\Delta t\to0$.

Bij een ECB verandert de grootte van $\vec v$ niet, maar wel zijn
richting. Daardoor is:

\[
\boxed{\vec a\neq\vec 0}.
\]


\FVDsubsection{De verandering van de snelheidsvector}

Beschouw twee nabije punten op de cirkel. De snelheidsvectoren
$\vec v_1$ en $\vec v_2$ hebben dezelfde grootte maar een iets andere
richting.

De verandering is:

\[
\Delta\vec v
=
\vec v_2-\vec v_1.
\]

Wanneer de twee punten steeds dichter bij elkaar komen, wijst
$\Delta\vec v$ steeds meer naar het middelpunt van de cirkel.

In de limiet geldt daarom:

\[
\boxed{
\vec a_c\ \text{wijst naar het middelpunt}.
}
\]

We noemen dit de \textbf{centripetale versnelling}.

``Centripetaal'' betekent letterlijk: naar het centrum gericht.


\begin{oefening}[Misconceptie: constante snelheid]{\oefU\ESS}

Een leerling zegt:

\begin{quote}
``De auto rijdt met constante snelheid door de ronde bocht. Volgens
$a=\Delta v/\Delta t$ is zijn versnelling dus nul.''
\end{quote}

Analyseer deze redenering.

\begin{question}
Welk onderscheid maakt de leerling niet?
\end{question}

\begin{question}
Welke eigenschap van $\vec v$ verandert?
\end{question}

\begin{question}
Teken twee snelheidsvectoren en construeer grafisch
$\Delta\vec v$.
\end{question}

\begin{question}
Formuleer een verbeterde redenering.
\end{question}

\end{oefening}


% ============================================================
% 8 AFLEIDING CENTRIPETALE VERSNELLING
% ============================================================

\FVDsection{De grootte van de centripetale versnelling}

We willen nu niet alleen de richting, maar ook de grootte van de
versnelling bepalen.

Tijdens een klein tijdsinterval $\Delta t$ verplaatst het voorwerp zich
over een kleine boog. Voor voldoende kleine $\Delta t$ geldt:

\[
\Delta s\approx v\Delta t.
\]

De positievector draait daarbij over een kleine hoek $\Delta\theta$.
Omdat

\[
\Delta s=r\Delta\theta,
\]

volgt:

\[
\Delta\theta
\approx
\frac{v\Delta t}{r}.
\]

De snelheidsvector draait over dezelfde hoek $\Delta\theta$. Voor een
kleine hoek is de grootte van de verandering van de snelheidsvector:

\[
|\Delta\vec v|
\approx
v\Delta\theta.
\]

Dus:

\[
|\Delta\vec v|
\approx
v\frac{v\Delta t}{r}
=
\frac{v^2}{r}\Delta t.
\]

Delen door $\Delta t$ en de limiet nemen geeft:

\[
\boxed{
a_c=\frac{v^2}{r}.
}
\]


\FVDsubsection{Andere vormen}

Met

\[
v=r\omega
\]

vinden we:

\[
a_c
=
\frac{(r\omega)^2}{r}
=
r\omega^2.
\]

Dus:

\[
\boxed{
a_c=r\omega^2.
}
\]

Met

\[
\omega=\frac{2\pi}{T}
\]

volgt:

\[
\boxed{
a_c=\frac{4\pi^2r}{T^2}.
}
\]

En omdat $f=1/T$:

\[
\boxed{
a_c=4\pi^2rf^2.
}
\]

Deze formules beschrijven hetzelfde fysische verband in verschillende
vormen.


\begin{eigenschap}[Eenparige cirkelbeweging]

Voor een ECB:

\[
\boxed{
v=r\omega
}
\]

en:

\[
\boxed{
a_c=\frac{v^2}{r}=r\omega^2
=\frac{4\pi^2r}{T^2}=4\pi^2rf^2.
}
\]

De vector $\vec a_c$ is steeds naar het middelpunt gericht.

\end{eigenschap}


\begin{opmerking}
Bij een ECB is de \emph{grootte} $a_c$ constant wanneer $r$ en $v$
constant zijn. De versnellingsvector $\vec a_c$ is echter niet constant,
want zijn richting verandert voortdurend.
\end{opmerking}


% ============================================================
% 9 VERBANDEN ANALYSEREN
% ============================================================

\FVDsection{Niet memoriseren, maar verbanden zien}

Het leerdoel vraagt meer dan formules invullen. We moeten kunnen
analyseren hoe een verandering van één grootheid andere grootheden
beïnvloedt.

Uit

\[
v=r\omega
\]

volgt bij constante $\omega$:

\[
v\sim r.
\]

Uit

\[
a_c=r\omega^2
\]

volgt bij constante $\omega$:

\[
a_c\sim r.
\]

Maar uit

\[
a_c=\frac{v^2}{r}
\]

volgt bij constante $r$:

\[
a_c\sim v^2.
\]

En uit

\[
a_c=\frac{4\pi^2r}{T^2}
\]

volgt bij constante $r$:

\[
a_c\sim\frac1{T^2}.
\]

De voorwaarde ``wat blijft constant?'' is dus essentieel.


\begin{oefening}[Wat gebeurt er als...?]{\oefU\ESS}

Beantwoord zonder eerst concrete getallen te kiezen.

\begin{question}
Bij constante straal verdubbelt de baansnelheid. Met welke factor
verandert $a_c$?
\end{question}

\begin{question}
Bij constante hoeksnelheid verdubbelt de straal. Met welke factor
veranderen $v$ en $a_c$?
\end{question}

\begin{question}
Bij constante straal halveert de periode. Met welke factor veranderen
$\omega$, $v$ en $a_c$?
\end{question}

\begin{question}
Bij constante baansnelheid verdubbelt de straal. Met welke factor
verandert $a_c$?
\end{question}

Motiveer elk antwoord vanuit een geschikt verband.

\end{oefening}


\begin{oefening}[De valkuil van ``de straal verdubbelt'']{\oefG}

Twee leerlingen bespreken wat er met de centripetale versnelling gebeurt
als de straal wordt verdubbeld.

Leerling A zegt:

\[
a_c=\frac{v^2}{r}
\quad\Rightarrow\quad
a_c\text{ halveert.}
\]

Leerling B zegt:

\[
a_c=r\omega^2
\quad\Rightarrow\quad
a_c\text{ verdubbelt.}
\]

\begin{question}
Kunnen beide redeneringen correct zijn?
\end{question}

\begin{question}
Welke impliciete voorwaarde gebruikt leerling A?
\end{question}

\begin{question}
Welke impliciete voorwaarde gebruikt leerling B?
\end{question}

\begin{question}
Formuleer een algemene les die je hieruit trekt bij het analyseren van
fysische formules.
\end{question}

\end{oefening}


% ============================================================
% 10 VOORBEELDEN
% ============================================================

\FVDsection{Toepassingen}

\begin{voorbeeld}[Een draaimolen]

Een zitje bevindt zich op $r=4,0\,\mathrm{m}$ van de draaias. De
draaimolen maakt één omwenteling in $6,0\,\mathrm{s}$.

\[
\omega
=
\frac{2\pi}{6,0}
\approx1,05\,\mathrm{rad/s}.
\]

De baansnelheid:

\[
v=r\omega
=
4,0\cdot1,05
\approx4,19\,\mathrm{m/s}.
\]

De centripetale versnelling:

\[
a_c
=
\frac{v^2}{r}
\approx
\frac{4,19^2}{4,0}
\approx4,39\,\mathrm{m/s^2}.
\]

De versnelling wijst op elk moment naar de draaias.

\end{voorbeeld}


\begin{oefening}[Hamerwerpen]{\oefB\ESS}

Een hamer beweegt vlak vóór het loslaten ongeveer in een cirkel met
straal $1,20\,\mathrm{m}$. De atleet maakt $5,0$ omwentelingen in
$2,0\,\mathrm{s}$.

\begin{question}
Bepaal de frequentie.
\end{question}

\begin{question}
Bepaal de periode.
\end{question}

\begin{question}
Bepaal de hoeksnelheid.
\end{question}

\begin{question}
Bepaal de baansnelheid.
\end{question}

\begin{question}
Bepaal de centripetale versnelling vlak vóór het loslaten.
\end{question}

\begin{question}
In welke richting beweegt de hamer onmiddellijk nadat hij wordt
losgelaten? Verklaar vanuit de snelheidsvector.
\end{question}

\end{oefening}


\begin{oefening}[Een bocht]{\oefU\ESS}

Een auto rijdt met constante snelheidsgrootte
$72\,\mathrm{km/h}$ door een cirkelvormige bocht met straal
$50\,\mathrm{m}$.

\begin{question}
Zet de snelheid om naar $\mathrm{m/s}$.
\end{question}

\begin{question}
Bepaal de centripetale versnelling.
\end{question}

\begin{question}
Schets $\vec v$ en $\vec a_c$ op drie verschillende punten van de
bocht.
\end{question}

\begin{question}
De bestuurder verdubbelt zijn snelheid terwijl de bochtstraal gelijk
blijft. Hoe verandert $a_c$?
\end{question}

\begin{question}
Waarom is de laatste conclusie verkeerskundig belangrijk?
\end{question}

\end{oefening}


% ============================================================
% 11 AARDE EN RUIMTE
% ============================================================

\FVDsection{Van de aarde tot satellieten}

Cirkelbeweging speelt een centrale rol in de astronomie.

Een punt op de evenaar beweegt door de aardrotatie ongeveer langs een
cirkel. Een satelliet in een bijna cirkelvormige baan beweegt eveneens
voortdurend met een versnelling naar het middelpunt van zijn baan.

Bij een satelliet is de relevante straal niet de \emph{hoogte} boven
het aardoppervlak, maar de afstand tot het middelpunt van de aarde:

\[
\boxed{
r=R_{\mathrm{aarde}}+h.
}
\]

Dat onderscheid is essentieel bij het mathematiseren.


\begin{oefening}[Een punt op de evenaar]{\oefU}

Neem voor de straal van de aarde:

\[
R_{\mathrm{aarde}}=6,37\times10^6\,\mathrm{m}
\]

en voor de rotatieperiode ongeveer $24,0\,\mathrm{h}$.

\begin{question}
Bepaal de hoeksnelheid van de aardrotatie.
\end{question}

\begin{question}
Bepaal de baansnelheid van een punt op de evenaar.
\end{question}

\begin{question}
Bepaal de centripetale versnelling van dat punt.
\end{question}

\begin{question}
Vergelijk deze versnelling met $g=9,81\,\mathrm{m/s^2}$.
\end{question}

\end{oefening}


\begin{oefening}[Satelliet: lees de straal correct]{\oefU\ESS}

Een satelliet bevindt zich $400\,\mathrm{km}$ boven het aardoppervlak.
Neem:

\[
R_{\mathrm{aarde}}=6,37\times10^6\,\mathrm{m}.
\]

De omlooptijd bedraagt $92,0\,\mathrm{min}$.

\begin{question}
Bepaal de straal van de cirkelbaan.
\end{question}

\begin{question}
Bepaal de hoeksnelheid.
\end{question}

\begin{question}
Bepaal de baansnelheid.
\end{question}

\begin{question}
Bepaal de centripetale versnelling.
\end{question}

\begin{question}
Waarom zou het fout zijn om $400\,\mathrm{km}$ als $r$ te gebruiken?
\end{question}

\end{oefening}


% ============================================================
% 12 PROBLEEMOPLOSSEN
% ============================================================

\FVDsection{Cirkelproblemen systematisch aanpakken}

Een formularium helpt alleen wanneer je eerst begrijpt welke grootheden
gegeven zijn en welke omstandigheden gelden.

Een bruikbare aanpak is:

\begin{enumerate}
  \item Maak een schets van de cirkelbaan.
  \item Duid het middelpunt en de straal aan.
  \item Noteer de gegeven grootheden met eenheden.
  \item Zet omwentelingen per minuut, uren, centimeters enzovoort eerst
        om naar geschikte SI-eenheden.
  \item Bepaal welke grootheden constant zijn.
  \item Kies het verband dat de gegeven en gezochte grootheden rechtstreeks
        met elkaar verbindt.
  \item Bereken en controleer de eenheid.
  \item Controleer de grootteorde.
  \item Interpreteer het antwoord fysisch, inclusief de richting wanneer
        het om een vector gaat.
\end{enumerate}


\begin{oefening}[Achterwaarts ontwerpen]{\oefG}

Een testopstelling moet een centripetale versnelling van
$3,0g$ produceren op een afstand van $0,75\,\mathrm{m}$ van de draaias.

Neem:

\[
g=9,81\,\mathrm{m/s^2}.
\]

\begin{question}
Bepaal de vereiste hoeksnelheid.
\end{question}

\begin{question}
Bepaal de vereiste frequentie.
\end{question}

\begin{question}
Druk het resultaat uit in omwentelingen per minuut.
\end{question}

\begin{question}
Bepaal de baansnelheid op $0,75\,\mathrm{m}$ van de as.
\end{question}

\begin{question}
Wat gebeurt er met $a_c$ voor een punt op $0,25\,\mathrm{m}$ van de
as als de hele opstelling met dezelfde $\omega$ draait?
\end{question}

\end{oefening}


% ============================================================
% 13 LABO
% ============================================================






% ============================================================
% 15 MISCONCEPTIES
% ============================================================

\FVDsection{Veelgemaakte denkfouten}

\begin{opmerking}[Misconceptie 1: ``eenparig'' betekent $a=0$]

Bij een ECB is de \emph{grootte} van de snelheid constant, maar de
richting verandert. Daarom is de versnelling niet nul.

\end{opmerking}

\begin{opmerking}[Misconceptie 2: $\vec a_c$ wijst in de bewegingsrichting]

De snelheidsvector is raaklijnig aan de baan. De centripetale
versnellingsvector staat daar loodrecht op en wijst naar het middelpunt.

\end{opmerking}

\begin{opmerking}[Misconceptie 3: centripetale versnelling is een extra soort beweging]

``Centripetaal'' beschrijft de richting van de versnelling. In het
volgende hoofdstuk onderzoeken we welke resulterende kracht deze
versnelling veroorzaakt.

\end{opmerking}

\begin{opmerking}[Misconceptie 4: $a_c$ is een constante vector]

Bij een ECB kan de grootte $a_c$ constant zijn, maar de vector
$\vec a_c$ verandert voortdurend van richting.

\end{opmerking}


\begin{oefening}[Vind de fout]{\oefU\ESS}

Beoordeel telkens de uitspraak en verbeter ze indien nodig.

\begin{question}
``Bij een ECB is de snelheid constant.''
\end{question}

\begin{question}
``Omdat $v$ constant is, is $a=0$.''
\end{question}

\begin{question}
``De centripetale versnelling is raaklijnig aan de cirkel.''
\end{question}

\begin{question}
``Als de baansnelheid verdubbelt bij dezelfde straal, verdubbelt ook
$a_c$.''
\end{question}

\begin{question}
``Twee punten op dezelfde stijve draaiende schijf hebben dezelfde
baansnelheid.''
\end{question}

\end{oefening}


% ============================================================
% 16 GEINTEGREERDE OEFENINGEN
% ============================================================

\FVDsection{Geïntegreerde oefeningen}

\begin{oefening}[Van gegevens naar model]{\oefU\ESS}

Een punt beweegt eenparig op een cirkel met straal
$0,80\,\mathrm{m}$. Het maakt $12$ omwentelingen in $15,0\,\mathrm{s}$.

\begin{question}
Bepaal $f$ en $T$.
\end{question}

\begin{question}
Bepaal $\omega$.
\end{question}

\begin{question}
Bepaal $v$.
\end{question}

\begin{question}
Bepaal $a_c$.
\end{question}

\begin{question}
Schets op één gekozen punt $\vec r$, $\vec v$ en $\vec a_c$.
\end{question}

\begin{question}
Hoe veranderen $v$ en $a_c$ als de frequentie verdubbelt terwijl de
straal gelijk blijft?
\end{question}

\end{oefening}


\begin{oefening}[Twee draaimolens vergelijken]{\oefU\ESS}

Draaimolen A heeft straal $r_A=3,0\,\mathrm{m}$ en periode
$T_A=6,0\,\mathrm{s}$.

Draaimolen B heeft straal $r_B=6,0\,\mathrm{m}$ en periode
$T_B=12,0\,\mathrm{s}$.

\begin{question}
Vergelijk hun hoeksnelheden.
\end{question}

\begin{question}
Vergelijk de baansnelheden aan de rand.
\end{question}

\begin{question}
Vergelijk de centripetale versnellingen aan de rand.
\end{question}

\begin{question}
Welke vergelijking zou je fout kunnen beantwoorden als je alleen naar
de grotere straal van B kijkt? Leg uit.
\end{question}

\end{oefening}


\begin{oefening}[Technisch ontwerp]{\oefG}

Een laboratoriumcentrifuge moet op een straal van
$0,12\,\mathrm{m}$ een centripetale versnelling van minstens
$500g$ bereiken.

\begin{question}
Bepaal de minimale hoeksnelheid.
\end{question}

\begin{question}
Bepaal de minimale frequentie.
\end{question}

\begin{question}
Druk de rotatiesnelheid uit in omwentelingen per minuut.
\end{question}

\begin{question}
Een monsterhouder wordt verplaatst naar $0,08\,\mathrm{m}$ van de as,
zonder de rotatiesnelheid te veranderen. Bepaal de nieuwe $a_c$.
\end{question}

\begin{question}
Welke parameter zou je moeten aanpassen om op de kleinere straal toch
$500g$ te behouden?
\end{question}

\end{oefening}


\begin{oefening}[Redeneren met een grafiek]{\oefG}

Een experiment levert bij constante straal een bijna rechte grafiek op
wanneer $a_c$ wordt uitgezet tegen $f^2$.

\begin{question}
Toon vanuit de theorie waarom dit verwacht wordt.
\end{question}

\begin{question}
Wat stelt de helling van de grafiek voor?
\end{question}

\begin{question}
Hoe kan je uit die helling de straal van de cirkelbaan bepalen?
\end{question}

\begin{question}
Wat zou een duidelijk niet-nul snijpunt met de verticale as kunnen
betekenen voor het experiment of het meetmodel?
\end{question}

\end{oefening}


% ============================================================
% 17 VERDIEPING WISKUNDIGE BESCHRIJVING
% ============================================================

\FVDsection{Verdieping: de cirkelbeweging als functie van de tijd}

Voor sterke leerlingen kunnen we de vectoriële beschrijving nog verder
uitwerken.

Kies een cirkel met middelpunt in de oorsprong en straal $r$. Als het
punt bij $t=0$ op de positieve $x$-as start en tegenwijzerzin beweegt,
dan is:

\[
\theta(t)=\omega t.
\]

De coördinaten zijn:

\[
\boxed{
x(t)=r\cos(\omega t)
}
\]

en:

\[
\boxed{
y(t)=r\sin(\omega t).
}
\]

De positievector is:

\[
\boxed{
\vec r(t)
=
r\cos(\omega t)\,\vec e_x
+
r\sin(\omega t)\,\vec e_y.
}
\]

Afleiden geeft:

\[
\vec v(t)
=
-r\omega\sin(\omega t)\,\vec e_x
+
r\omega\cos(\omega t)\,\vec e_y.
\]

De grootte is:

\[
|\vec v|
=
r\omega
\sqrt{\sin^2(\omega t)+\cos^2(\omega t)}
=
r\omega.
\]

Nogmaals afleiden:

\[
\vec a(t)
=
-r\omega^2\cos(\omega t)\,\vec e_x
-
r\omega^2\sin(\omega t)\,\vec e_y.
\]

Dus:

\[
\boxed{
\vec a(t)=-\omega^2\vec r(t).
}
\]

Dit resultaat maakt onmiddellijk zichtbaar dat de versnelling tegengesteld
gericht is aan de positievector: naar het middelpunt.

Bovendien:

\[
|\vec a|
=
\omega^2r
=
a_c.
\]


\begin{oefening}[Vectorieel bewijzen]{\oefG}

Vertrek van:

\[
\vec r(t)
=
r\cos(\omega t)\,\vec e_x
+
r\sin(\omega t)\,\vec e_y.
\]

\begin{question}
Bepaal $\vec v(t)$ door af te leiden.
\end{question}

\begin{question}
Toon aan dat $|\vec v|=r\omega$.
\end{question}

\begin{question}
Bepaal $\vec a(t)$.
\end{question}

\begin{question}
Toon aan dat $\vec a=-\omega^2\vec r$.
\end{question}

\begin{question}
Leg uit waarom dit zowel de grootte als de richting van de
centripetale versnelling bewijst.
\end{question}

\end{oefening}


% ============================================================
% 18 SAMENVATTING
% ============================================================

\FVDsection{Wat hebben we opgebouwd?}

Voor een periodieke beweging:

\[
\boxed{
f=\frac1T
}
\]

Voor een ECB:

\[
\boxed{
\omega=\frac{2\pi}{T}=2\pi f
}
\]

De baansnelheid:

\[
\boxed{
v=r\omega=\frac{2\pi r}{T}
}
\]

De snelheidsvector is raaklijnig aan de cirkel:

\[
\boxed{
\vec v\perp\vec r.
}
\]

Hoewel $|\vec v|$ constant is, verandert $\vec v$ voortdurend van
richting. Daarom is er een versnelling.

De centripetale versnelling wijst naar het middelpunt:

\[
\boxed{
\vec a_c\ \text{naar het centrum}
}
\]

en haar grootte is:

\[
\boxed{
a_c=\frac{v^2}{r}
}
\]

of equivalent:

\[
\boxed{textb
a_c=r\omega^2
=\frac{4\pi^2r}{T^2}
=4\pi^2rf^2.
}
\]

Bij het analyseren van deze verbanden moet je altijd expliciet aangeven
welke grootheden constant worden gehouden.





% ============================================================
% 20 VOORUITBLIK
% ============================================================

\FVDsection{Vooruitblik: wat houdt het voorwerp op de cirkel?}

We weten nu wat nodig is voor een cirkelbeweging:

\[
\vec a_c
\]

moet naar het middelpunt gericht zijn.

Maar we hebben nog niet verklaard \emph{waardoor} die versnelling
ontstaat.

Bij een auto in een bocht, een steen aan een touw, een satelliet rond
de aarde en een passagier in een draaimolen zijn verschillende
fysische krachten actief.

In het volgende hoofdstuk verschuift de vraag daarom van

\[
\boxed{\text{Hoe beweegt het voorwerp?}}
\]

naar

\[
\boxed{\text{Waarom verandert zijn beweging?}}
\]

Daarvoor hebben we de wetten van Newton nodig.

\end{document}
